\documentclass[12pt]{article}
\usepackage[utf8]{inputenc}
\usepackage[T1]{fontenc}
\usepackage{hyperref}
\title{Ushojo Language Resources for K-12 Curriculum Development}
\author{Compiled by Claude Sonnet 5 (claude-sonnet-5) for Andr\'as Kornai}
\date{}
\begin{document}
\maketitle
\section{Introduction}
Ushojo (also Ushoji; ISO 639-3 ush) is an Indo-Aryan language spoken in the Swat and Kohistan districts of Khyber Pakhtunkhwa, Pakistan, by about 2,000 people. Search found only reference and encyclopedic pages; no primer, textbook, dictionary or recordings.
\section{Primary Educational Resources}
None found in search results.
\section{Linguistic and Reference Materials}
\begin{itemize}
\item Wikipedia entry (describes a Perso-Arabic script adapted from Torwali and Shina alphabets; this is not a standardized orthography): \url{https://en.wikipedia.org/wiki/Ushoji_language}
\item Endangered Languages Project entry, citing Decker (1992), ``Ushojo'' in \emph{Languages of Kohistan}, as the main descriptive source; it lists no standard orthography and no recent resources: \url{https://endangeredlanguages.com/node/96632}
\end{itemize}
\section{Religious and Community Resources}
None found in search results.
\section{Digital and Audio Resources}
None found in search results.
\section{Curriculum Development Context}
The language is not known to be taught in schools. Pashto exerts pressure on it, though some sources report children still learning it at home. Any curriculum would require new documentation and orthography development with the community; Decker (1992) is the starting reference.
\end{document}